% GENERATED FILE. Edit canonical lessons, then run npm run generate:books.
% canonical-source-hash: fnv1a64:c94dc6b48b01bcb6
% canonical-lessons: AR-C274-tabia2, AR-C274-hibr, AR-C274-sabbura, AR-C274-qamus, AR-C274-malaff

\chapter{Printer, Ink, Blackboard, Dictionary, File}
\label{ch:ar-printer-ink-blackboard-dictionary-file}
\hlchaptermodality{274}

\begin{chapteropening}
\textbf{By the end of this chapter:} \emph{I can say a printer, ink, a blackboard, a dictionary and a file.}

Complete the last lesson of chapter 274: I can say a printer, ink, a blackboard, a dictionary and a file.
\end{chapteropening}

\section[\texorpdfstring{\ar{طابعة} (ṭābiʿa)}{طابعة (ṭābiʿa)}]{\ar{طابعة} (ṭābiʿa) — a printer}

\label{lesson:AR-C274-tabia2}



\begin{quote}
Before the new one: say the Arabic for flour, then the Arabic for spices.
\end{quote}

\subsection*{You'll want to know: \ar{طابعة}}
\textbf{\ar{طابعة}} — \emph{ṭābiʿa} — \textquotedblleft{}a printer\textquotedblright{}.

\begin{cousinweb}
\textbf{\ar{طابعة}} (\emph{ṭābiʿa}) is built on the native Arabic root \emph{ṭ-b-ʿ}, \textquotedblleft{}to seal, to stamp\textquotedblright{} (Lane). \emph{Ṭābiʿa} is the feminine active participle, \textquotedblleft{}a thing that prints\textquotedblright{}.
\end{cousinweb}

\begin{grammarlens}[title={What you've built}]
One more word for this chapter.
\end{grammarlens}

\subsection*{Your turn}
\begin{itemize}
\raggedright
  \item \emph{Say it:} \emph{ṭābiʿa}
  \item \emph{Say it:} \emph{ṭābiʿa}, once more
  \item \emph{Say it:} say \emph{ṭābiʿa}
  \item \emph{Recall:} say the Arabic for flour, then the Arabic for spices, then say \emph{ṭābiʿa} again
\end{itemize}

\begin{tcolorbox}[breakable,colback=teal!4,colframe=teal!35!black,title={Before you move on}]
What does \ar{طابعة} mean? (\textquotedblleft{}A printer\textquotedblright{}.) Say it once more.
\end{tcolorbox}
\section[\texorpdfstring{\ar{حبر} (ḥibr)}{حبر (ḥibr)}]{\ar{حبر} (ḥibr) — ink}

\label{lesson:AR-C274-hibr}



\begin{quote}
Before the new one: say the Arabic for spices, then the Arabic for a printer.
\end{quote}

\subsection*{You'll want to know: \ar{حبر}}
\textbf{\ar{حبر}} — \emph{ḥibr} — \textquotedblleft{}ink\textquotedblright{}.

\begin{cousinweb}
\textbf{\ar{حبر}} (\emph{ḥibr}) is built on the native Arabic root \emph{ḥ-b-r}, \textquotedblleft{}ink; to adorn\textquotedblright{} (Lane). \emph{Ḥibr} is ink.
\end{cousinweb}

\begin{grammarlens}[title={What you've built}]
One more word for this chapter.
\end{grammarlens}

\subsection*{Your turn}
\begin{itemize}
\raggedright
  \item \emph{Say it:} \emph{ḥibr}
  \item \emph{Say it:} \emph{ḥibr}, once more
  \item \emph{Say it:} say \emph{ḥibr}
  \item \emph{Recall:} say the Arabic for spices, then the Arabic for a printer, then say \emph{ḥibr} again
\end{itemize}

\begin{tcolorbox}[breakable,colback=teal!4,colframe=teal!35!black,title={Before you move on}]
What does \ar{حبر} mean? (\textquotedblleft{}Ink\textquotedblright{}.) Say it once more.
\end{tcolorbox}
\section[\texorpdfstring{\ar{سبورة} (sabbūra)}{سبورة (sabbūra)}]{\ar{سبورة} (sabbūra) — a blackboard}

\label{lesson:AR-C274-sabbura}



\begin{quote}
Before the new one: say the Arabic for a printer, then the Arabic for ink.
\end{quote}

\subsection*{You'll want to know: \ar{سبورة}}
\textbf{\ar{سبورة}} — \emph{sabbūra} — \textquotedblleft{}a blackboard\textquotedblright{}.

\begin{cousinweb}
\textbf{\ar{سبورة}} (\emph{sabbūra}) is a modern word for a blackboard; it has no entry in Lane, and its origin is not settled.
\end{cousinweb}

\begin{grammarlens}[title={What you've built}]
One more word for this chapter.
\end{grammarlens}

\subsection*{Your turn}
\begin{itemize}
\raggedright
  \item \emph{Say it:} \emph{sabbūra}
  \item \emph{Say it:} \emph{sabbūra}, once more
  \item \emph{Say it:} say \emph{sabbūra}
  \item \emph{Recall:} say the Arabic for a printer, then the Arabic for ink, then say \emph{sabbūra} again
\end{itemize}

\begin{tcolorbox}[breakable,colback=teal!4,colframe=teal!35!black,title={Before you move on}]
What does \ar{سبورة} mean? (\textquotedblleft{}A blackboard\textquotedblright{}.) Say it once more.
\end{tcolorbox}
\section[\texorpdfstring{\ar{قاموس} (qāmūs)}{قاموس (qāmūs)}]{\ar{قاموس} (qāmūs) — a dictionary}

\label{lesson:AR-C274-qamus}



\begin{quote}
Before the new one: say the Arabic for ink, then the Arabic for a blackboard.
\end{quote}

\subsection*{You'll want to know: \ar{قاموس}}
\textbf{\ar{قاموس}} — \emph{qāmūs} — \textquotedblleft{}a dictionary\textquotedblright{}.

\begin{cousinweb}
\textbf{\ar{قاموس}} (\emph{qāmūs}) is usually traced to Greek \emph{ōkeanos}, \textquotedblleft{}ocean\textquotedblright{} (EtymDB). The dictionary sense comes from al-Fīrūzābādī's \emph{al-Qāmūs al-Muḥīṭ}, \textquotedblleft{}the surrounding ocean\textquotedblright{}.
\end{cousinweb}

\begin{grammarlens}[title={What you've built}]
One more word for this chapter.
\end{grammarlens}

\subsection*{Your turn}
\begin{itemize}
\raggedright
  \item \emph{Say it:} \emph{qāmūs}
  \item \emph{Say it:} \emph{qāmūs}, once more
  \item \emph{Say it:} say \emph{qāmūs}
  \item \emph{Recall:} say the Arabic for ink, then the Arabic for a blackboard, then say \emph{qāmūs} again
\end{itemize}

\begin{tcolorbox}[breakable,colback=teal!4,colframe=teal!35!black,title={Before you move on}]
What does \ar{قاموس} mean? (\textquotedblleft{}A dictionary\textquotedblright{}.) Say it once more.
\end{tcolorbox}
\section[\texorpdfstring{\ar{ملف} (malaff)}{ملف (malaff)}]{\ar{ملف} (malaff) — a file, a folder}

\label{lesson:AR-C274-malaff}



\begin{quote}
Before the new one: say the Arabic for a blackboard, then the Arabic for a dictionary.
\end{quote}

\subsection*{You'll want to know: \ar{ملف}}
\textbf{\ar{ملف}} — \emph{malaff} — \textquotedblleft{}a file, a folder\textquotedblright{}.

\begin{cousinweb}
\textbf{\ar{ملف}} (\emph{malaff}) is built on the native Arabic root \emph{l-f-f}, \textquotedblleft{}to wrap, to roll up\textquotedblright{} (Lane). \emph{Malaff} is a thing to wrap papers in: a file.
\end{cousinweb}

\begin{grammarlens}[title={What you've built}]
One more word for this chapter.
\end{grammarlens}

\subsection*{Your turn}
\begin{itemize}
\raggedright
  \item \emph{Say it:} \emph{malaff}
  \item \emph{Say it:} \emph{malaff}, once more
  \item \emph{Say it:} say \emph{malaff}
  \item \emph{Recall:} say the Arabic for a blackboard, then the Arabic for a dictionary, then say \emph{malaff} again
\end{itemize}

\begin{tcolorbox}[breakable,colback=teal!4,colframe=teal!35!black,title={Before you move on}]
What does \ar{ملف} mean? (\textquotedblleft{}A file, a folder\textquotedblright{}.) Say it once more.
\end{tcolorbox}
